% RESUMO - PT
\begin{resumo}
  \vspace{-15pt}

  O registro eletrônico em saúde concentra grande parte da informação clínica em
  texto livre, no qual a relação entre conceitos, em especial a negação de
  achados, é determinante para a correta interpretação. Este trabalho investiga
  a extração de relações entre entidades clínicas em português sobre o corpus
  SemClinBr, considerando o espaço de rótulos \texttt{negation\_of},
  \texttt{associated\_with} e \texttt{no\_relation}. Como etapa inicial do
  desenvolvimento do RECLin-PT, um modelo aprimorado de extração de relações
  clínicas em português, conduz-se um experimento controlado para responder a uma
  questão de projeto: o pré-treinamento clínico importa para essa tarefa? Propõe-se
  uma cadeia experimental reprodutível que vai do texto bruto às métricas finais,
  composta por um \emph{parser} determinístico do formato XML do corpus, análise
  exploratória das relações, particionamento 80/10/10 em nível de documento,
  estratificado e congelado para evitar vazamento de dados, e geração
  controlada de pares candidatos. Sobre essa infraestrutura, dois \emph{encoders}
  para o português são ajustados variando apenas o \emph{checkpoint} de
  pré-treino: um de domínio clínico (BioBERTpt) e um de domínio geral (BERTimbau).
  No conjunto de teste congelado, os dois empatam estatisticamente no F1 de
  \texttt{negation\_of} na semente de referência (0,677 contra 0,694;
  \emph{bootstrap} pareado IC95\% $[-0{,}053; +0{,}018]$, $p=0{,}345$), mas esse
  veredito não se sustenta nas duas sementes testadas. Na semente 43 a diferença
  \emph{inverte de sinal} e passa a favorecer o \emph{encoder} clínico de forma
  significativa (IC95\% $[+0{,}016; +0{,}090]$). A variação entre sementes de um
  mesmo modelo supera a diferença entre os modelos, de modo que nem o veredito
  nem sua direção são estáveis \emph{nessa classe}, ao passo que, no padrão
  global de acertos, o teste de McNemar aponta para o \emph{encoder} geral nas
  duas sementes. Esses resultados indicam que o RECLin-PT pode ser
  construído sobre o \emph{encoder} geral, que acompanha o clínico em macro-F1
  médio com cerca de 60\% dos parâmetros e é cerca de dez vezes mais estável
  entre reinicializações.

  \vspace{\onelineskip}

  \noindent\textbf{Palavras-chave}: \palavraschaveemlinha
\end{resumo}

% RESUMO - EN
\begin{resumo}[Abstract]
  \vspace{-15pt}

  \begin{otherlanguage*}{english}

  Electronic health records store much clinical information as free text, in which
  the relationship between concepts, particularly the negation of findings,
  is decisive for correct interpretation. This work investigates relation
  extraction between clinical entities in Portuguese over the SemClinBr corpus,
  considering the label space \texttt{negation\_of}, \texttt{associated\_with} and
  \texttt{no\_relation}. As the initial stage of developing RECLin-PT, an improved
  clinical relation extraction model for Portuguese, we run a controlled
  experiment to answer a design question: does clinical pre-training matter for
  this task? We propose a reproducible experimental pipeline that goes from raw
  text to final metrics, comprising a deterministic parser for the corpus XML
  format, exploratory analysis of relations, a document-level, stratified and
  frozen 80/10/10 split to prevent data leakage, and controlled candidate-pair
  generation. On this infrastructure, two Portuguese encoders are fine-tuned
  changing only the pre-trained checkpoint: a clinical-domain one (BioBERTpt) and
  a general-domain one (BERTimbau). On the frozen test set, the two are
  statistically tied on the \texttt{negation\_of} F1 under the reference seed
  (0.677 vs.\ 0.694; paired bootstrap 95\% CI $[-0.053, +0.018]$, $p=0.345$), but
  this verdict does not hold across the two seeds we tested. Under seed 43 the
  difference \emph{flips sign} and significantly favours the clinical encoder
  (95\% CI $[+0.016, +0.090]$). Seed-to-seed variation of a single model exceeds
  the gap between the models, so neither the verdict nor its direction is stable
  \emph{on that class}, whereas on the overall pattern of correct predictions
  McNemar's test points to the general encoder under both seeds.
  These results indicate that RECLin-PT can be built on the general encoder,
  which matches the clinical one in mean macro-F1 with about 60\% of the
  parameters and is roughly ten times more stable across restarts.

  \vspace{\onelineskip}

  \noindent\textbf{Keywords}: \inlinekeywords
\end{otherlanguage*}
\end{resumo}
